% Einheit 2 von 3 – Binomische Formeln (bank/binomische-formeln/e2.jsonl, zusammenbau v0.1)
%% TODO Zweigzeile Teil 1: was hier gelernt wird (Satz in Schülersprache aus dem Einheitstitel) – Einheit 2
\einheitenkopf[e2]{Einheit 2 von 3 · Binomische Formeln}
\zweigzeile{ab Kl. 8 · P10}
\setcounter{aufgabe}{12}

%% TODO Titel: Ich-kann-Satz fehlt in der Bank (Platzhalter: Kettenname) – Nr. 13
%% TODO Anweisung: ein Satz über der ersten Teilaufgabe fehlt in der Bank – Nr. 13
\begin{aufgabe}{Gleiche Klammer zweimal?}
\begin{teile}
\teil $(x + 9)^2$ – Quadrat einer Klammer? Wenn ja: schreib es als Klammer mal Klammer. \janein
\end{teile}
\end{aufgabe}

%% TODO Titel: Ich-kann-Satz fehlt in der Bank (Platzhalter: Kettenname) – Nr. 14
%% TODO Anweisung: ein Satz über der ersten Teilaufgabe fehlt in der Bank – Nr. 14
\begin{aufgabe}{Binomische Formeln}
\begin{teile}
\teil Welche binomische Formel passt, und was ist $a$, was ist $b$ in der Formel? $(x + 14)^2$ \leerfeld Formel; a = \leerfeld , b = \leerfeld
\end{teile}
%% TODO Darstellung für schwach fehlt in der Bank (binomische-formeln-e2-k2-s1-v1; Abschnitt „Für schwache Schüler“)
\swz{$\text{Multipliziere aus: $(x + 4)^2$}$}{}
%% TODO Darstellung für schwach fehlt in der Bank (binomische-formeln-e2-k2-s1-v2; Abschnitt „Für schwache Schüler“)
\swz{$\text{Multipliziere aus: $(x + 9)^2$}$}{}
%% TODO Darstellung für schwach fehlt in der Bank (binomische-formeln-e2-k2-s1-v3; Abschnitt „Für schwache Schüler“)
\swz{$\text{Multipliziere aus: $(x + 1)^2$}$}{}
%% TODO Darstellung für schwach fehlt in der Bank (binomische-formeln-e2-k2-s1-v4; Abschnitt „Für schwache Schüler“)
\swz{$\text{Multipliziere aus: $(x + 8)^2$}$}{}
%% TODO Darstellung für schwach fehlt in der Bank (binomische-formeln-e2-k2-s1-v5; Abschnitt „Für schwache Schüler“)
\swz{$\text{Multipliziere aus: $(x + 11)^2$}$}{}
\swfrage{Was bleibt gleich, was ändert sich?}
%% TODO Darstellung für schwach fehlt in der Bank (binomische-formeln-e2-k2-s2-v1; Abschnitt „Für schwache Schüler“)
\swz{$\text{Multipliziere aus: $(x - 9)^2$}$}{}
%% TODO Darstellung für schwach fehlt in der Bank (binomische-formeln-e2-k2-s3-v1; Abschnitt „Für schwache Schüler“)
\swz{$\text{Multipliziere aus: $(x + 8) \cdot (x - 8)$}$}{}
%% TODO Darstellung für schwach fehlt in der Bank (binomische-formeln-e2-k2-s4-v1; Abschnitt „Für schwache Schüler“)
\swz{$\text{Welche Formel passt? Multipliziere aus: $(x + 9) \cdot (x - 9)$}$}{}
%% TODO Darstellung für schwach fehlt in der Bank (binomische-formeln-e2-k2-s5-v1; Abschnitt „Für schwache Schüler“)
\swz{$\text{Multipliziere aus: $(3x + 1)^2$}$}{}
%% TODO Darstellung für schwach fehlt in der Bank (binomische-formeln-e2-k2-s6-v1; Abschnitt „Für schwache Schüler“)
\swz{$\text{Multipliziere aus: $(x + 3y)^2$}$}{}
%% TODO Darstellung für schwach fehlt in der Bank (binomische-formeln-e2-k2-s7-v1; Abschnitt „Für schwache Schüler“)
\swz{$\text{Multipliziere aus: $2 \cdot (x + 4)^2$}$}{}
%% TODO Darstellung für schwach fehlt in der Bank (binomische-formeln-e2-k2-s8-v1; Abschnitt „Für schwache Schüler“)
\swz{$\text{Multipliziere aus: $-(x + 9)^2$}$}{}
%% TODO Darstellung für schwach fehlt in der Bank (binomische-formeln-e2-k2-s9-v1; Abschnitt „Für schwache Schüler“)
\swz{$\text{Multipliziere aus und fasse zusammen: $(x + 1)^2 - 6$}$}{}
\swz[3]{Die Parabel $p$ hat die Gleichung $y = (x + 4)^2 - 7$. Weise nach, dass $y = x^2 + 8x + 9$ dieselbe Parabel beschreibt. \hfill (P10 2017 OS)}{}
\end{aufgabe}

%% TODO Titel: Ich-kann-Satz fehlt in der Bank (Platzhalter: Kettenname) – Nr. 15
%% TODO Anweisung: ein Satz über der ersten Teilaufgabe fehlt in der Bank – Nr. 15
\begin{aufgabe}{Kopfrechnen mit der Formel (Vorrat)}
%% TODO Darstellung für schwach fehlt in der Bank (binomische-formeln-e2-k3-s1-v1; Abschnitt „Für schwache Schüler“)
\swz{$41^2$ – mit einer binomischen Formel im Kopf?}{}
\end{aufgabe}

%% TODO Titel: Ich-kann-Satz fehlt in der Bank (Platzhalter: Kettenname) – Nr. 16
%% TODO Anweisung: ein Satz über der ersten Teilaufgabe fehlt in der Bank – Nr. 16
\begin{aufgabe}{Binomische Formeln – Fehler finden, Begründen}
\swz[3]{Finde den Fehler und rechne richtig: \rechnung{(y + 9)^2 &= y^2 + 81}}{}
\swfrage{Tim meint: $(x + 4)^2 = x^2 + 16$. Prüfe mit $x = 2$ und begründe, warum das nicht stimmt.}
\end{aufgabe}

\uebersichtskasten{Erste binomische Formel: (a + b)² = a² + 2ab + b² \quad (x + 3)² = x² + 6x + 9 \\ Zweite binomische Formel: (a − b)² = a² − 2ab + b² \quad (x − 3)² = x² − 6x + 9 \\ Dritte binomische Formel: (a + b)·(a − b) = a² − b² \quad (x + 3)·(x − 3) = x² − 9 \\ Das Mittelglied 2ab nie vergessen – (a + b)² ist nicht a² + b². Zahlenprobe: (2 + 3)² = 25, aber 2² + 3² = 13. \\ Mit Vorzahl ist a die ganze Vorzahl mit x. \quad (2x + 3)² = 4x² + 12x + 9}
